% =====================================================================
% Section 5 — Curve25519 in Microcode
% Draft for TCHES submission.
%
% Companion to Section 4 (Keccak): the second worked example that
% instantiates the general methodology of Section 3. Where Keccak stressed
% control flow and register residency, Curve25519 stresses the dense
% multiply-accumulate-carry chain of reduced field arithmetic. All
% Curve25519-specific detail lives here. Numbers and design rationale are
% taken from the field-op patches (asm_op_curve25519.c / _mul.c), our
% fastest X25519 implementation (full_curve25519_inline2.c, the
% register-chained all-inline-assembly ladder = the "ucode-inline"
% contender), the same-ladder isolation harness (full_curve25519.c), and
% the recorded matrix scoreboard (RESULTS.md).
%
% Assumes the Background has recalled X25519 / the Montgomery ladder and
% the field GF(2^255-19); we restate only what the mapping needs.
% =====================================================================

\section{Curve25519 in Microcode}
\label{sec:curve25519}

The second worked example instantiates the methodology of
\Cref{sec:ucode-layer} on the field arithmetic that underlies X25519.
Whereas Keccak (\Cref{sec:keccak}) was realised as a single large,
control-flow-heavy patch looped in place, a Curve25519 field operation has
the opposite character. It is a small, straight-line, multiplication-dominated
kernel that is invoked many hundreds of times within a single scalar
multiplication. The primitive that we encode in microcode is consequently the
field multiplication and squaring modulo $2^{255}-19$, and the question we
address is whether a microcode patch can evaluate these operations more
quickly than the best available native code, namely compiler output, an
automatically generated and formally verified backend, a superoptimiser, and
expert hand-written assembly.

\paragraph{The structural advantage.}
In this setting the advantage of the layer is not register residency, since
five limbs occupy only a small fraction of the register file. The
performance-critical path of a reduced multiplication is a sum of partial
products that is accumulated into a small number of running totals and then
propagated as carries, and on native x86 the limiting factor is the carry
rather than the arithmetic. Because the architecture exposes a single carry
flag, multi-word carry propagation is forced to serialise through the
\texttt{ADC} instruction, and a two-ALU in-order core cannot overlap that
dependent chain with the multiplications that feed it. The microcode layer
addresses this limitation in two ways, both of which follow from the
properties established in \Cref{sec:findings}. A single triad fuses a
multiplication, an accumulation, and a carry capture into one issue bundle,
and the per-register \texttt{SETCC} flag domains allow several carry chains to
progress concurrently without contending for one architectural flag. This
concurrency of carry propagation is the field-arithmetic counterpart of the
register residency that Keccak exploited, in that it is work the instruction
set architecture is unable to express.

% ---------------------------------------------------------------------
\subsection{Representation and Operations}
\label{sec:c25519-rep}

We adopt the unsaturated radix-$2^{51}$ representation, in which a field
element is stored as five 64-bit limbs that each carry 51 significant bits.
This is the representation emitted by fiat-crypto and used by the
Bernstein--Schwabe \texttt{amd64-51} implementation. Multiplication is the
$5\times5$ schoolbook of 25 partial products, and the reduction implied by
$2^{255}\equiv 19$ is folded directly into the cross terms so that no separate
reduction pass is required, giving
\[
c_i \;=\; \sum_{j\le i} a_{\,i-j}\,b_j \;+\; 19\!\!\sum_{j>i} a_{\,5+i-j}\,b_j .
\]
Squaring exploits the symmetry $a_ib_j=a_ja_i$ to reduce the 25 products to
15, and the doublings $2a_i$ and folds $19a_i$ that this requires are
precomputed by the caller and supplied in registers. The carries that arise
during accumulation are short, because each limb is reduced to 51 bits and the
remaining 13-bit overflow, which together with the limb width fills a 64-bit
word, is shifted into the following limb. The brevity and independence of
these carries are what make the representation well suited to the layer, a
point to which we return in \Cref{sec:c25519-perf}, since it is exactly where
the $4\times64$ saturated representation behaves differently.

% ---------------------------------------------------------------------
\subsection{The Patch: Fused Multiply--Accumulate and Parallel Carry Chains}
\label{sec:c25519-patch}

A complete field multiplication or squaring is carried out by a single
execution of one hooked macro-instruction, an event we term a firing. The
native ladder triggers the operation with that one instruction, whereupon the
corresponding patch evaluates the entire reduced product in registers, which is
to say all of the partial products together with the carry propagation and the
reduction modulo $2^{255}-19$, and returns control to the native side with the
result in place. There is in consequence neither any per-limb looping nor any
repeated round-trip between native code and microcode within a single field
operation. The two patches reside in the patch RAM simultaneously, the
multiplication occupying 66
triads from \texttt{U7c00} and the squaring 42 triads from \texttt{U7d08}, and
they are therefore reached through two distinct match registers that redirect
two different macro-instructions, in the manner described in
\Cref{sec:triads-packing}. The multiplication is hooked to \texttt{vmwrite},
whose opcode is \texttt{0x0cd8}, while the squaring is hooked to
\texttt{vmread}, whose opcode is \texttt{0x0618}, so that the native code fires
the multiplication by issuing a \texttt{vmwrite} and the squaring by issuing a
\texttt{vmread}. Using a separate entry instruction for each patch is what
allows both to be installed at once, since a single hooked instruction can be
redirected to only one entry point. The caller passes the five limbs, together
with the precomputed $19b_j$ folds in the case of multiplication, in fixed
registers. A short
preparation stage then materialises the remaining folds in place, computing
$19b_1,\dots,19b_4$ with three multiplications packed into one triad by means
of the slot-write rule described in \Cref{sec:findings}, after which the kernel
traverses the limbs and emits for each of them a chain of fused
multiply--accumulate triads of the form shown in \Cref{lst:c25519-mac}.

\begin{lstlisting}[language=C,caption={The fused multiply--accumulate--carry
core of the squaring kernel for a single limb. The register \texttt{TMP0}
holds the low accumulator, \texttt{R8} the high accumulator, \texttt{TMP9} the
running carry sum, and \texttt{TMP15} the transient \texttt{SETCC} carry. A
multiplication, an accumulation, and a carry capture share one
triad.},label={lst:c25519-mac},basicstyle=\ttfamily\footnotesize]
{ ADD(TMP0,TMP0,RDI), SETCC(TMP15,TMP0), ADD(R8,R8,RCX),  NOP_SEQW } // acc lo; carry; acc hi
{ ZEROEXT(TMP9,TMP15),MUL(RCX,RBX,R13),  ADD(TMP0,TMP0,R13),NOP_SEQW } // next product while accumulating
\end{lstlisting}

\paragraph{Three concurrent accumulators.}
The decisive feature of the kernel is that the sum of partial products is
maintained as three independent chains, none of which touches the single
architectural flag. The low halves of the products accumulate into
\texttt{TMP0}, the high halves accumulate into \texttt{R8}, and the carries
produced by the individual accumulations are summed into \texttt{TMP9}. Each
chain captures its own carry by means of a \texttt{SETCC} into a per-register
flag domain, as established in \Cref{sec:findings}, with the consequence that
the carry capture occupies a spare slot of the multiplication's triad rather
than serialising behind it. The transition between limbs is handled by a
single \texttt{OR} that seeds the next limb's low accumulator from the two
shifted pieces of the previous limb's carry, namely the high part of the
accumulator shifted right by 51 bits combined with the carry sum shifted left
by 13 bits. A native multiplication bound by \texttt{ADC} cannot reproduce
this scheme, because it would be obliged to thread every carry through one
flag.

\paragraph{Authoring and verification.}
In contrast to the Keccak permutation, the Curve25519 field patches are
hand-scheduled rather than generator-emitted. The kernels are small enough, at
a few tens of triads, to lay out by hand, and the dense and irregular packing
of multiplications and carries leaves little for a generic packer to improve.
They are nevertheless held to the discipline described in
\Cref{sec:triads-packing}. Each patch is verified against the fiat-crypto
reference, which is also the backend that the superoptimiser targets, over a
large number of random inputs, and the assembled X25519 is checked against the
known-answer vectors of RFC~7748 before any performance claim is made.

% ---------------------------------------------------------------------
\subsection{Triad Budget}
\label{sec:c25519-budget}

The squaring kernel compiles to 42 triads and the multiplication kernel to 66
triads, and both therefore lie comfortably within the 128-triad capacity of
the patch RAM. This is the qualitative respect in which Curve25519 differs
from Keccak, because the patch budget is not the binding constraint and the
governing design variable is not how to make the operation fit but how densely
it can be packed. Every multiply--accumulate--carry sequence that is folded
into a single triad shortens both the patch and its critical path, and the
optimisation history of these kernels consists largely of such merges. The
multiplication, for example, was reduced from 73 to 66 triads solely by
packing the preparation stage's multiplications three to a triad and by
drawing the per-product carry capture into the free slots of the
multiply--accumulate triads.

% ---------------------------------------------------------------------
\subsection{Integration into X25519}
\label{sec:c25519-integration}

A field multiplication or squaring is invoked several hundred times in the
course of a single scalar multiplication, which comprises a 255-step
Montgomery ladder with several field operations per step followed by a Fermat
inversion that is itself dominated by squarings. Because each invocation is a
single hooked firing, a \texttt{vmwrite} for a multiplication and a
\texttt{vmread} for a squaring, the per-firing dispatch overhead of
approximately four cycles established in \Cref{sec:keccak} is amortised against
the cost of the field operation.

Our complete implementation, which is the fastest that we constructed and which
we denote the inlined ladder, realises the entire scalar multiplication as
inline assembly built around these two patches rather than as C that calls them
through function wrappers. The fifteen field elements that a ladder step
manipulates are held together in a single stack-resident state structure at
fixed offsets, and the assembly addresses them through a base register that is
pinned to that structure for the duration of the step. The microcode firings
are emitted inline within one large assembly block per ladder step, with the
consequence that the compiler spills registers only around the block as a whole
rather than at every field operation. Within that block, consecutive operations
are chained through registers wherever the dataflow permits, in that the output
limbs of one operation are passed directly as the inputs of the next and the
reload that would otherwise begin the second operation is omitted; a single
ladder step contains five such chains, of which one example is the addition
forming $A=x_2+z_2$ feeding directly into the squaring $AA=A^2$. The Fermat
inversion is realised in the same manner, as one inlined assembly block in which
the long sequence of squarings is chained through registers, so that only the
first squaring loads from memory and only the last stores its result. This
arrangement extends the principle of \Cref{sec:budget}, namely that data
movement is the quantity to be minimised, from the interior of a single field
operation to the boundaries between successive operations, and it is the reason
that the inlined ladder is faster than the otherwise identical implementation
that invokes the field operations through function calls.

In order to measure the contribution of the field operation in isolation, which
the inlined ladder by its construction entangles with the integration, we
additionally conduct a controlled comparison in which the ladder, the inversion
chain, and the conditional swap are held fixed and only the field backend is
varied. This same-ladder comparison is conducted against four alternatives,
namely hand-written \texttt{\_\_uint128\_t} C, the automatically generated C
produced by fiat-crypto, the Goldmont-tuned assembly produced by the CryptOpt
superoptimiser, and the hand-written \texttt{amd64-51} assembly of Bernstein and
Schwabe. We perform the comparison both within our own ladder and within the
\texttt{amd64-51} ladder, so that the microcode backend is evaluated against
expert assembly on the framework of the assembly's own authors.

% ---------------------------------------------------------------------
\subsection{Performance}
\label{sec:c25519-perf}

\paragraph{Methodology.}
We follow the measurement discipline of \Cref{sec:keccak}, timing at the base
frequency within a single process and sweeping over 24
compiler-and-optimisation configurations while retaining the best result for
each contender, in the same manner as the SUPERCOP autotuner. Since the
comparison is conducted at the level of a complete X25519, the headline figure
is the geometric mean of the per-configuration ratios, which collapses the
grid to a single number for each backend.

\paragraph{The field-operation result.}
When the ladder is held fixed, the microcode field operations are the fastest
backend that we measured. Within our own ladder the microcode is faster than
hand-written C by a factor of $1.24$, faster than fiat-crypto by a factor of
$1.20$, and faster than the CryptOpt superoptimiser by a factor of $1.22$,
each expressed as a geometric mean over the 24 configurations and reported in
\Cref{tab:c25519-sameladder}. The microcode therefore outperforms optimised C,
a formally verified automatically generated backend, and a state-of-the-art
superoptimiser, with every component other than the field operation held
constant. When the microcode is substituted for the hand-written assembly in
the \texttt{amd64-51} framework of Bernstein and Schwabe, it remains faster by
a factor of $1.07$, so that the advantage persists even against expert
hand-tuned x86 assembly. The aggregate improvement is driven by the squaring,
which in isolation runs in 37 cycles and is therefore $26\%$ faster than
fiat, while the multiplication runs in 58 cycles and is $3\%$ faster; since a
ladder is dominated by squarings, the squaring advantage governs the overall
result.

\begin{table}[t]
\centering
\caption{Same-ladder field-operation comparison, reported as the geometric
mean of the per-configuration minimum-cycle ratios for a complete X25519,
where a smaller value indicates faster execution and the field backend is the
only component that varies within each block.}
\label{tab:c25519-sameladder}
\small
\begin{tabular}{@{}llr@{}}
\toprule
ladder & field-operation backend & ratio vs.\ microcode \\
\midrule
ours        & hand-written \texttt{\_\_uint128\_t} C & $1.24\times$ \\
ours        & fiat-crypto (autogenerated C)          & $1.20\times$ \\
ours        & CryptOpt (superoptimised asm)          & $1.22\times$ \\
ours        & \textbf{microcode}                     & $1.00\times$ \\[2pt]
amd64-51    & Bernstein--Schwabe hand asm            & $1.07\times$ \\
amd64-51    & \textbf{microcode}                     & $1.00\times$ \\
\bottomrule
\end{tabular}
\end{table}

\paragraph{The integration result.}
The inlined ladder of \Cref{sec:c25519-integration} is the fastest of our
implementations, completing X25519 in $300{,}197$ cycles, as against $304{,}897$
cycles for the otherwise identical implementation that invokes the same
microcode field operations through function calls, as reported in
\Cref{tab:c25519-e2e}. The register chaining and the inlining of the firings
therefore account for the difference of approximately $1.5\%$ between the two,
a margin consistent with the expectation that store-to-load forwarding hides
most, though not all, of the inter-operation memory traffic.

\paragraph{The limits of the result.}
Considered end to end, however, the fastest X25519 on Goldmont is not our own.
That position is held by the Bernstein--Schwabe \texttt{amd64-64}
implementation, which uses the $4\times64$ saturated representation with radix
$2^{64}$ and completes in 271k cycles, ahead of our fastest implementation, the
inlined microcode ladder, at approximately 300k cycles, as reported in
\Cref{tab:c25519-e2e}. A
more troubling observation for the layer is that when we port the microcode
field operations to that same $4\times64$ representation they regress by
roughly a factor of two relative to the native assembly. The explanation lies
precisely in the carry structure discussed in \Cref{sec:c25519-rep}, because a
saturated $4\times64$ multiplication requires long, full-width carry chains
across 64-bit limbs, which is the pattern that our flag model handles least
well. Since \texttt{ADC} reads the architectural flag domain that is frozen at
patch entry, as noted in \Cref{sec:findings}, a four-deep saturated carry must
be emulated with serial \texttt{SETCC} operations, and the advantage conferred
by parallel carry propagation is thereby lost. The unsaturated $5\times51$
representation is favourable to the microcode for the converse reason, namely
that its 13-bit carries remain short and mutually independent. The
representation that is most favourable to native assembly is consequently the
one at which the layer is weakest, and the reverse holds as well.

\begin{table}[t]
\centering
\caption{End-to-end X25519 performance, given as the best result for each
contender over 24 configurations in minimum cycles. The field-operation
advantage does not overcome the representation-level advantage of
\texttt{amd64-64}, and the microcode $4\times64$ port is where the layer's
single-flag carry model is exposed.}
\label{tab:c25519-e2e}
\small
\begin{tabular}{@{}llr@{}}
\toprule
contender & representation and backend & min (cyc) \\
\midrule
amd64-64 / asm        & $4\times64$ saturated, hand asm        & 271\,314 \\
ours / ucode-inline   & $5\times51$, microcode (inlined)       & 300\,197 \\
ours / ucode          & $5\times51$, microcode                 & 304\,897 \\
amd64-51 / ucode      & $5\times51$, microcode (B--S ladder)   & 332\,323 \\
donna\_c64            & $5\times51$, portable C                & 335\,736 \\
ours / fiat           & $5\times51$, fiat-crypto C             & 354\,874 \\
amd64-51 / asm        & $5\times51$, hand asm                  & 355\,975 \\
ours / cryptopt       & $5\times51$, CryptOpt asm              & 379\,941 \\
amd64-64 / ucode      & $4\times64$ saturated, microcode       & 510\,750 \\
\bottomrule
\end{tabular}
\end{table}

\paragraph{Summary.}
At the level at which the paper makes its claim, namely the field operation
with the surrounding ladder, inversion, and conditional swap held fixed, the
microcode is the fastest Curve25519 field backend that we measured on
Goldmont, surpassing optimised C, fiat-crypto, the CryptOpt superoptimiser,
and expert hand-written assembly. The result has the same character as that of
Keccak in \Cref{sec:keccak}, in that it derives from a structural
per-operation advantage that the instruction set architecture cannot reach,
which in this instance is the fused multiply--accumulate together with the
parallel \texttt{SETCC} carry chains rather than register residency. As with
Keccak, the boundary of the result is stated honestly. The end-to-end
advantage belongs to a different representation, the $4\times64$ saturated
form, and that representation is exactly where the one structural weakness of
the layer, namely long carry chains under a single-flag emulation model,
becomes apparent.
